% !TeX root = ../教材题库_全部条目_转写工作稿.tex
% Source: XJ-B1-C06-EX-P252-Y0348
% 第6章 / 6.4 用样本估计总体 / 教材 p.252、253 / PDF 第259、260页
% OCR未校对：公式与文字均可能有误；未完成数学TeX转写；禁止直接发给学生
% 原题及相邻原解截图（核对用）：../题目原图/XJ-B1-C06-EX-P252-Y0348-01.webp, ../题目原图/XJ-B1-C06-EX-P252-Y0348-02.webp
% 本文件只在转写完成后作为题干使用；OCR 条目尚不是可靠题干。
\begin{bankitem}{6.4 例题 例题11}{第6章 / 6.4 用样本估计总体 / 教材 p.252、253 / PDF 第259、260页}{XJ-B1-C06-EX-P252-Y0348}{OCR未校对：公式与文字均可能有误；未完成数学TeX转写；禁止直接发给学生}
\begin{ocrdraft}
[教材p.252 / PDF第259页]
\par
0：，二]为满足居民基本用电需求，并提高能源的利用效率，实现绿色发展，市
\par
政府计划对全市居民用电标准按年采用三阶式递增电价收费：75％的用户在最低电
\par
价一档，20％的用户用电量超出一阶电价的临界值而未超过二阶电价的临界值，超
\par
过一阶临界值的用电量按二阶电价缴费，5％的用户用电量超过二阶电价的临界值，
\par
超过二阶临界值的用电量按三阶电价缴费．为此，当地电力公司调查了200户居民6
\par
月份的用电量（单位：
\par
．h)，并排序如下：
\par
8
\par
22
\par
31
\par
42
\par
48
\par
9
\par
50
\par
51
\par
56
\par
57
\par
57
\par
60
\par
61
\par
61
\par
61
\par
62
\par
62
\par
66
\par
67
\par
70
\par
70
\par
72
\par
72
\par
74
\par
76
\par
77
\par
77
\par
78
\par
78
\par
80
\par
82
\par
82
\par
82
\par
84
\par
84
\par
88
\par
88
\par
89
\par
90
\par
91
\par
93
\par
93
\par
94
\par
96
\par
96
\par
96
\par
98
\par
98
\par
98
\par
99
\par
100
\par
100
\par
100
\par
101
\par
101
\par
101
\par
105
\par
106
\par
106
\par
106
\par
107
\par
107
\par
107
\par
107
\par
108
\par
108
\par
109
\par
109
\par
110
\par
]彐0
\par
110
\par
111
\par
112
\par
]彐3
\par
113
\par
]彐4
\par
115
\par
116
\par
118
\par
120
\par
120
\par
120
\par
121
\par
123
\par
124
\par
127
\par
127
\par
127
\par
130
\par
130
\par
130
\par
131
\par
131
\par
132
\par
132
\par
132
\par
133
\par
133
\par
134
\par
134
\par
13
\par
135
\par
135
\par
135
\par
135
\par
136
\par
137
\par
137
\par
138
\par
139
\par
139
\par
]过0
\par
141
\par
142
\par
14
\par
146
\par
146
\par
147
\par
148
\par
149
\par
151
\par
152
\par
巧4
\par
156
\par
159
\par
160
\par
162
\par
163
\par
163
\par
164
\par
165
\par
167
\par
169
\par
170
\par
170
\par
172
\par
174
\par
17
\par
177
\par
178
\par
178
\par
180
\par
182
\par
182
\par
187
\par
189
\par
191
\par
191
\par
192
\par
194
\par
194
\par
200
\par
201
\par
201
\par
202
\par
203
\par
203
\par
206
\par
208
\par
212
\par
213
\par
214
\par
216
\par
223
\par
224
\par
237
\par
247
\par
250
\par
250
\par
251
\par
253
\par
254
\par
258
\par
260
\par
265
\par
274
\par
274
\par
283
\par
288
\par
289
\par
304
\par
319
\par
320
\par
324
\par
339
\par
462
\par
498
\par
530
\par
542
\par
626
\par
阶梯电价的临界点如何确定？
\par
分析电力公司选取6月份进行调查，是因为6月用电量在一年12个月中处于
\par
中等偏上（可能需要空调用电），一年的用电量近似等于这个月的用电量乘以12．阶
\par
梯电价临界值的确定依赖于总体分布，我们要用样本数据的信息确定第75百分位数
\par
以及第95百分位数．
\par
解样本容量7冖200．
\par
为了使75％的用户以最低一档（第一档）电价缴费，需确定P7\ensuremath{\cdot}
\par
因为c一”\ensuremath{\times}75％一150是整数，
\par
150一+一工151
\par
178十178
\par
所以P75
\par
=17\&
\par
2
\par
2
\par
[教材p.253 / PDF第260页]
\par
这说明6月份约有75％的用户用电量不超过178kW．h.假设每年用电量不超过
\par
178\ensuremath{\times}12=2136（kW．h）按第一档电价缴费，则大约75％的用户将按第一档缴费．
\par
因为c一刀\ensuremath{\times}95％一190是整数，
\par
190十工“289十304
\par
所以P9
\par
=296，5．
\par
2
\par
2
\par
这说明6月份约有95％的用户用电量不超过296，5kW\ensuremath{\cdot}h,也就是说该月用电
\par
量超过296．5kW\ensuremath{\cdot}h的用户只有5％．若每年用电量超过2136kW．h，不超过
\par
296．5\ensuremath{\times}12=3558（kW\ensuremath{\cdot}h)，则其中的2136kW．h按第一档电价缴费，超过部分
\par
按第二档电价缴费．
\par
若每年用电量超过3558kW．h，则超过部分按第三档电价缴费．
\end{ocrdraft}
\end{bankitem}
